\documentclass[11pt,a4paper]{article}
\usepackage{amsmath, amssymb, amsthm}
\usepackage[authoryear]{natbib}
\usepackage[margin=1in]{geometry}
\usepackage{booktabs}
\usepackage[hidelinks]{hyperref}
\usepackage{setspace}

\setlength{\parindent}{0pt}
\setlength{\parskip}{6pt}
\onehalfspacing

%opening
\title{Felt Passage Without a Spotlight: Discrete Neural Frames and One-Way Dependency in the Block Universe}
\author{Tyce Hanna}

\begin{document}

\maketitle

\begin{abstract}

Consciousness ignites when the metabolic threshold is met and goes dark when that threshold isn't met (e.g., anesthesia). Between those poles, the brain parses continuous sensory streams into sub-second processing windows, synthesizing the data into a single perceptual frame. That frame is not a spectator of raw stimulus but rather is what the stimulus becomes after integration, and it can only point upstream. It contains encoding from prior frames and none of what has not been encoded. Continuity is a conviction already written into those states.

This paper treats felt time as a property of that discrete, one-way process. Duration is perceived change relative to a species-dependent perceptual ceiling and a logged baseline. When integration falls toward zero, experience ceases. Boredom, flow, and acute stress are read as changes in how full that ceiling is and how cleanly data are integrated. Critical flicker-fusion thresholds are used as a rough comparative proxy for frame rate across individuals and species. Dopamine and norepinephrine enter only as interpretative modulators of integration but do not generate temporal order. A formal biological equation is proposed to relate change, baseline, and ceiling to the felt passage of time.

The philosophical setup is kept narrow. If past, present, and future are equally real in a static spacetime block, standard theories still incorporate external variables to attempt to explain why time feels like it flows and why memory only runs backward: entropy as an independent arrow, or a moving spotlight that then needs a speed and a second time to measure it, leading to hyper time and empty regression. The claim here does not deny entropy, thermodynamic asymmetry, or spacetime physics. It locates the smooth, one-way experience of time within neurology via discrete perceptual cycles and one-way encoding. From inside the static 4D block, that is what the passage of time feels like.

\end{abstract}

\section{Introduction}

Einstein wrote that for ``believing physicists'' the separation of past, present, and future has ``only the meaning of an illusion, albeit a tenacious one'' \citep{einstein1955besso}. Instead, past, present, and future are all equally and simultaneously real, coexisting at the same moment within a frozen block of spacetime, also called the block universe. If correct, this raises an obvious question: If the universe is a static block and all moments exist simultaneously, how do we experience the smooth passage of time?

This paper proposes an alternative. Rather than introducing external variables, the answer lies within neurology itself. The mechanics of consciousness—and how our brain integrates data—provide a simple and inherent explanation for the fluid experience of time and why we can only remember the past.

This account would be weakened if felt duration endured when integration dropped to near zero if percepts could be shown to depend on downstream rather than upstream content, or if a continuous, non-chunked stream of consciousness were required to explain ordinary temporal experience.

Existing temporal direction models are reviewed first. Next, premises about the spacetime block, the world line, and thresholded consciousness are stated. The mechanism in focus is the discrete frames of perception, the one-way dependency of a percept, the continuity of those frames, and why time does not flash by. Finally, the paper covers ordinary experience, modulators across species, and thought experiments. Limits and implications close the paper.

\section{Existing models of temporal directionality}

The standard models for time perception and directionality are classified in three main groups: B theory, C theory, and moving spotlight theory, each of which is incomplete.

It’s important to differentiate between the three arrows of time: physical, thermodynamic, and psychological. The physical time arrow is the parameter that appears in spacetime structure. The thermodynamic arrow is the familiar increase of entropy described by the second law. The psychological arrow is the felt order of experience and why memory is restricted to the past. This paper addresses the third. 

Standard B theory correctly acknowledges that the past, present, and future all exist equally within a static block, with no special “now,” but fails to explain why an observer feels a fluid direction of time. The theory waves off directionality as an illusion, which does no explanatory work. Dismissing temporal flow as an illusion produced by the being embedded within the block is no more informative than pointing at a photograph and saying the motion it depicts is an illusion produced by the printing process \citep{mctaggart1908unreality}.

More developed B theory models attempt to explain directionality with the second law of thermodynamics, suggesting the psychological arrow of time results from increased entropy, which tends to only go in one direction for complex systems. While this is stronger than bare B theory, it incorporates a separate external explanatory principle (entropy), which may not be necessary at all. 

Standard C theory models state that while events are ordered, there is no inherent directionality to these events. Instead, temporal asymmetry results from entropy only increasing in one direction \citep{Dummett1960-DUMADO}.

Moving spotlight theory agrees the past, present, and future exist equally and attempts to explain ``now'' and the arrow of time through a present moment spotlight that moves across the static block, illuminating each moment in turn \citep{skow2009relativity}.

However, this moving light presents a challenge. If the spotlight moves, it must have a speed, which means it must also have another time dimension in which to measure it. This creates an infinite regress known as hyper time.

Standard B theory suffers from a similar problem but tends to wave it off as an illusion. However, this is an empty claim. Even though B theory acknowledges the 4D block is completely motionless, the intuitive feeling of moving forward through time still requires some explantion.

A pattern emerges here: ``time is an illusion,'' ``one second per second,'' and ``the block is a perfect mirror of the 3D world'' on their own are empty statements. While entropy attempts to address directionality, it leaves the question of how a static structure gives rise to a fluid experience. None of the theories discussed here provide an inherent explanation for the fluid illusion of directionality experienced by conscious beings. While entropy is an attempt to explain psychological asymmetry, it may not be required once neural dependency is in place.

\section{Premises}

When reading this, it’s important to keep in mind that the self is not a fixed thing, rather the unique neurological lens through which an observer perceives reality.

In this paper, I treat the 4D spacetime block as a perfect static mirror of the 3D and all its contents, rather than treating the block as some alternate reality. Although it is tempting to claim this alone explains how we perceive the flow of time, it, by itself is an empty statement.

The premises explored in this paper can be labelled as follows:

\begin{itemize}
	\item P1. The 4D spacetime block contains the same events and neural states as the 3D world. It’s not a second, distinct reality.
	\item P2. At any event on an observer’s world line, only the neural record encoded at the event is available. There is no access to contents that have not been encoded.
	\item P3. Consciousness is present when the relevant neural conditions are met and absent in anesthesia, dreamless sleep, and coma.
\end{itemize}

The mirror premise, properly interpreted, makes things much easier to understand. Nobody requires an explanation for why we can’t see inside a stranger’s home, or why we can’t see a month into the future. The answer is simple: neurologically, the mind has zero record of these contents, therefore, they cannot be perceived. We all understand, neurologically, the brain only has the contents of the past up to point x.

Imagine a scanner moving across a paper. There is a set amount of data per its own activity, again not implying movement as within the block there is a permanent static record of the scanner’s entire history, but the point remains, a scanner halfway through the paper, only has read x amount of the page at x point.

An observer’s relationship inherits this unremarkable principle. At any given point within a world line, the observer only has access to what has happened, not because of entropy or some external mechanism, but because he simply lacks neurological access to what hasn't happened.

The concept of “frame” is used in two related ways. The primary use is a neural processing window, a discrete integrated percept like that described by time-slice and perceptual-cycle research. The secondary use is analogical where each such percept is a static snapshot at a location on the world line.

\section{Mechanism}

\subsection{Discrete frames}

Neuroscience suggests we perceive reality in discrete chunks (or frames) via a two-stage model of perception. The importance of this ``chunking'' is to convert chaotic multimodal sensory data into a coherent cognitive unit, organizing what would otherwise be sensory chaos \citep{herzog2016}. Without this chunking mechanism, the conscious experience would collapse into an incoherent, overwhelming blur of sensory static.

Modern neuroscience suggests a dual mechanism where the brain collects sensory stimuli unconsciously over a time slice, then the brain integrates that data, processes it, and delivers it as a single discrete frame. Alpha and theta brain waves act like a refresh rate or shutter speed of human consciousness. The brain samples the environment rhythmically at about 10-13 frames per second (alpha rhythm) \citep{vanrullen2016}. This means what we perceive as a smooth continuous ``movie'' or experience is actually comprised of a rapid sequence of discrete information packages. Gamma band brain oscillations bind the disparate sensory inputs across different parts of the brain into a singular, discrete unit of conscious perception \citep{doesburg2009}.

These perceptual chunks act as static conscious snapshots within the spacetime block. Each of these individual conscious snapshots has a record of the frames prior but not after. For the same reason we can’t know what the interior of a house looks like if we’ve never been inside.

\subsection{Dependency}

A percept is the neurologically processed result of a given stimulus. It cannot exist prior or independent of what it results from and has no coherent meaning independent of what came before it. A prior model described the percept as a separate state that perceived stimulus after some interval, but in that case, the claim quietly implies a gap between a raw unprocessed stimulus and a later state that perceived it, inviting an infinite regress of perceiving states. A percept is not some separate state but rather is what integrated data is labeled. The stimulus and percept are not two separate cognitive states, but are instead the same fixed fact within the block, described from two angles: 1) raw unprocessed data and 2) what that data becomes following cognitive integration \citep{jepson1993percept}.

For example, a sculpture depends on and is logically downstream from the marble it was carved from, even though there is no gap between them to bridge once the structure exists. The statue is always dependent on the marble instead of the other way around. The marble does not depend on the sculpture—that dependency only runs one way. A percept within the block can only reference upstream percepts. There are only fixed static coordinates within the block, which the contents encode.

The percept can reference data as far back as memory will allow, from the earliest static frames in which neurology encoded a long-term memory. The incremental changes within these frames are what results in the continuity we experience.

\subsection{Why time feels continuous}

The reasonable follow-up question is to ask why these separate static conscious frames produce the continuous smooth experience we know instead of just being separate, static content. This is likely the result of viewing consciousness and the self through a rigid inaccurate lens, which is similar to another question: Why am I me? This is a common human fallacy we make because we assume we had to be chosen, but there was no cosmic waiting room we were selected from. You wouldn't ask why is this rock this specific rock, so why would you ask, “Why am I me?” It’s no different, you just are, from the outside its very simple, but feels complex from the inside.

This applies to frames within the static block as well. They are not separate, fixed individuals with some unique soul; they are the localized phenomenon that is consciousness. A frame does not compare itself to some alternate, disconnected version of itself. It’s not like it can look around and see its separate structure relative to the other frames. A frame is simply a very small unit of data only capable of perceiving the small frozen unit of data within itself, which again, is downstream of prior frames, like a conscious snapshot, not because it was selected or moved, but because that conviction is what the neurological state within a frame consists of.

The frame prior to, say, standing up and closing a door is convinced of its continuity and that in x time, the door will be closed. The hypothetical frame after getting up and closing a door is completely convinced that it just closed the door because all the frames within that period of getting up and shutting it provide the microscopic, frame-by-frame changes necessary to create a continuous experience. These frames contain all the neurologic data of each moment from getting up to close the door and closing it, even though nothing actually moved. The very last frame after shutting the door includes the belief that the subject walked over to the door and closed it.

Walking is another familiar case of felt motion. When we walk, it feels like a continuous self is producing motion through spacetime. Despite how convincing it feels, there is no motion occurring. Instead, each frame is a complete static percept, including physical positioning and memory. Within the block, every point on the worldline contains neurologically integrated stimulus.

\subsection{Why time does not flash by}

The question how does time not flash by is a fallacy, based on the idea that our perception of time is some magical external thing, but the real-world data refutes this clearly, when perception is at or close to zero (e.g., dreamless sleep, anesthesia, coma), our perception of time ceases to exist. This is because the internal feeling of time passing is purely a result of data perception, since your worldline contains all the data within your lifetime and the neurological states that perceived the data, it is fundamentally impossible for time to just ``flash'' by. Flashing by and time not existing are two sides of the same coin.

Within the block, our brain and its frozen units integrate and react to stimulus whether internal or external. The block contains a frozen replica of every moment, the brain perceived that stimulus, and each of those conscious snapshots has a fraction of that integrated data. Together, they build a seamless illusion that time passed. Not to say that clocks are responsible for time perception, but that integrated data aligns with that measurement in the frozen neurological state within the block.

\section{The mechanism in ordinary experience}

Our perception of time can be reduced to two pillars, internal and external. The external pillar includes outside stimulus such as sound, touch, or sight, while the internal pillar consists of internal wandering and bodily signals (e.g., heartbeat, hunger, etc.). However, these two pillars aren't independent, co-equal channels but combine to create a synthesized experience. Without external stimulus, or remaining internal stimulus, the mind does not have any content to generate its own contrast either, nothing to remember, anticipate, or reflect on.

However, the brain does not simply maintain a set value for each area of perception but rather allocates energy to available areas. In some cases, internal activity is increased when external is limited. For example, the brains of people with blindness allocate energy to other areas, such as hearing or tactile sensation \citep{khalaf2025shared}. Sensory deprivation also supports this. While the subject may be deprived of new external stimulus, they have a lifetime of stimulus experience to draw from mentally.

This mechanism is not limited to extreme or clinical cases; it is active in many common human experiences. Consider a slow workday. The environment may offer little external stimulus and instead of the brain integrating data relative to the low stimulus, it begins generating its own, such as watching the clock, thinking about when work ends, checking the temperature, and so on. As a result, waves of stress hormones are released, which act like a knob cranking up sensory stimulus data overall \citep{henckens2009}. As a result, temporal responses are amplified, and time feels like it slows down.

Conversely, when stimulus integration drops to near zero, such as during sleep, anesthesia, or a coma, time seems to cease to exist. For example, a student in a boring, under-stimulating class starts playing a video game. As boredom and distress have already caused a net increase in sensory data, playing a game (a predictable activity) releases dopamine, which in turn assists in allocating attention away from the boredom. This allows the brain to filter out incoming sensory data and stop the release of stress hormones, which causes a net decrease in the integrated data. Suddenly, the class that felt excruciatingly boring and slow now begins moving at a normal pace \citep{qiu2018rapid}. A similar effect can be observed when checking a phone while waiting in line or listening to the radio while sitting in a car.

A more unusual example is the geometric tunnels commonly reported during intense psychedelic experiences. Under this sensory and cognitive disruption, a highly stimulated brain seeks additional data and appears to generate its own in the form of complex, changing geometric tunnels. They provide a sense of velocity, change, depth, color, and rich complexity to deliver large amounts of stimulus to the brain. The experience seems to be intensified when the eyes are closed, providing a dark black canvas for the geometry by limiting external data \citep{kluver1928mescal}.

This neurological mechanism may explain why a wide variety of substances and experiences produce similar geometry. However, ths is speculative and not confirmed.

\section{Modulators and comparative proxies}

\subsection{The fixed rule and the perceptual ceiling}

Total perceived change relative to a species perceptual ceiling and past perceptual baseline is the most precise formulation of the mechanism. Subjective time is not generated by some external force or some cross species baseline, it is produced by how full that unique, species-dependent ceiling is relative to its logged prior baseline. There is no universal ceiling. Each species just needs its perceptual ceiling filled at a proportional level.

How this ceiling is applied varies greatly among species. A housefly, for example, perceives at 200hz to 250hz, while a human only perceives at around 60hz. However, perception quality varies drastically between the species. Human vision is analogous to a 4k UltraHD TV while a housefly’s is like a small, grainy screen. Despite the disparity, the housefly doesn't need a huge perceptual ceiling. Even though its perception quality is lower, it can evade threats effectively.

How different species perceive time can be calculated with deviation. For example, take a 1980s beat up truck, and a Bugatti. Whether you’re cruising at 300 mph or 60mph, you don't feel any change, you only feel change when  slowing down or speeding up. Time perception is no different. The only reference a species has is prior memories compared to recent moments.

To calculate the entirety of time perception across species, you must calculate the perceptual ceiling, which consists of five primary data input variables: somatic, external sensory, mental thought, predictive processing, and subcortical. To calculate this across species, we need significantly more data than what we currently do \citep{healy2013metabolic}.

The visual system is the best way to roughly measure this across typical mammals. The field of view (FOV), the ocular resolution, and the critical flicker-fusion frequency (CFF) are all direct representations of a species perceptual ceiling. This is similar to how the processing capacity of a PC determines how a video game runs. Nature does not evolve ocular capabilities beyond what a species can neurologically handle, just like you can roughly infer the specs of a PC based on the frames per second (FPS) requirements of a video game.

However, to accurately gauge the perceptual ceiling, we need data on how species value various inputs. For example, while both bats and humans are mammals, bats heavily favor hearing compared to other senses. To calculate deviation, we must calculate the memory window scale and the depth of prior memory data across a given amount of time. For example, if you sit in a dark shed for five seconds and step back out, you don't feel as though time got noticeably faster. However, if you sit in the shed for a month, once you step out, you will experience cognitive shock as the weight of that memory window was much higher and the memory comparison index registers a massive change relative to past rate.

Time perception likely requires advanced neural architecture to maintain the barrier between the self and the stimulus. Without this boundary, an organism cannot compare the present moment to the past and has no concept of time because there is no individual to perceive time. Worms, for example, may not have that boundary and therefore may lack temporal perception.

This may explain why users of psychedelics report a point when the self is completely dissolved. The boundary between the self and stimulus evaporates, creating the perception that time has expanded to infinity (or the opposite).

Although visual metrics are the easiest proxy to measure, calculating deviation would require advanced neuroimaging to test the various sensory inputs across different species. As far as measuring deviation is concerned, the best metric we have is behavioral data. Behavioral data, such as reaction time and excitability, represents a species' perceptual variation.

Flies, for example, show very low variance in behavioral data, which aligns with their incentives as a species: they must always be able to evade threats at any given moment to survive and reproduce. Certain species, such as hummingbirds, sit on the opposite side of the spectrum. These birds show very high behavioral variance depending on the environment and context \citep{gonzalezgomez2011flexibility}.

Because the hummingbird's survival depends on catching very small and fast insects, its excitability and reaction time raise drastically. But when night comes, they enter a state of torpor. Because their baseline metabolism burns through such high amounts of fuel and they lack significant fat stores like other birds, their metabolic rate drops drastically, and along with it, so do their behavioral data rates \citep{shankar2022heterothermic}.

\subsection{Flicker fusion rate across species}

If the internal rate of perception of time is governed by processing speed that varies through individual biology, it should be measurable and vary widely across species based on biological incentives. The best current proxy for this is CFF, the rate at which a flickering light stops being perceived as individual flashes and appears as a continuous glow. CFF is one of the few proxies for perceptual ceiling that has a wide range of established data across species.

A broad study found values from under 1 hz to several hundred hz in fast flying insects. Across species, CFF correlates with how the organism lives and survives. Vulnerable, faster moving species cluster towards the high end while less vulnerable, slower species tend to cluster at the lower end \citep{lafitte2022}. For example, human CFF is four times slower than the housefly, which is why they are so hard to swat \citep{haarlem2024speed}. A human hand swatting at it represents a modest speed compared to the rate at which its eyes (and brain) update \citep{coile1989}.

This framework treats CFF as a rough metric for the perceptual ceiling as it has established data across species and requires significant processing power. Calculating the perceptual ceiling requires significantly more data than just CFF, as sensory input varies wildly across species, even within the same group. CFF is a visual proxy. It is strongest for species that depend heavily on vision and weakest for species in which auditory or olfactory senses dominate. Bats, for example, primarily use smell and hearing (echolocation) to survive while humans are more dependent on vision. Ocular, auditory, olfactory, somatic, and internal mental activity all contribute to the ceiling’s total sum of sensory data.

\subsection{Dopamine as the active modulator of data tallying}

While the previous section addressed the ceiling itself, this section addresses how species allocate and integrate energy and data based on neurological factors. The primary neurological factor is dopamine, although the neurological architecture as a whole is the main factor. While the precise role is not fully established, dopamine seems to govern how the brain allocates its energy and integrates data, whether its effective or not is context dependent, based on the demands of the current moment.

Dopamine's established role in attention motivation and reward has been well documented, and dopaminergic dysfunction is a recognized role in ADHD \citep{volkow2011}. Dopamine may also play a role in how the neurological architecture within the brain develops \citep{silberstein2016}.

Dopamine plays a significant role as an active modulator of how much stimulus is tallied cognitively, how that stimulus is integrated into a coherent trackable signal, and how brains allocate energy. Low dopamine does not necessarily mean less is being perceived, but what is perceived is being integrated less effectively and functions closer to noise than a usable integrated signal, producing the exact symptoms of attention deficit hyperactivity disorder (ADHD): disorganization, forgetfulness, and difficulty sustaining attention.

This framing also makes sense of an apparent contradiction in how ADHD presents itself. The same person who struggles to track unstimulating, repetitive tasks can sometimes sustain hours of unbroken focus. If dopamine governed only the total amount of perceived stimulus, this would look inconsistent. However, a more accurate frame may be that dopamine governs the total sum of integrated data. For example, an ADHD student may hear the teacher but struggle to track and integrate what has been said. This is also why, for an ADHD person, hours can pass in what felt like half the time while engaging in stimulating activities (e.g., playing video games). Dopamine signaling allocates focus to the highly predictable, low entropy video game and tunes out the rest of the world because the ceiling is less full than it otherwise would be \citep{qiu2018rapid}.

In contrast, the same person in an under-stimulating classroom may feel as though time is dragging because dopamine is lower and the tasks are not absorbing. The person then checks the clock, retreats into their mind, notices their clothing against their skin, feels the uncomfortable seat, and releases stress hormones like cortisol—all of which dramatically increase the internal data stream \citep{henckens2009}.

Norepinephrine is another important neurotransmitter involved in time perception. As a highly relevant factor in ADHD, Norepinephrine is a key mechanism of action in prescription ADHD meds and may play a more complimentary and narrowing role in the process. Where dopamine in large part modulates allocation and integration effectiveness, norepinephrine appears to govern how wide the attentional/perceptual net is cast \citep{ghosh2024locus}.

In cases of life-threatening emergencies (i.e., high entropy, low familiarity situations), norepinephrine and other stress hormones narrow the size of the attentional “net” to focus and allow the brain to perceive an acute threat in detail. In this narrow frame, the brain calculates at a higher rate than it otherwise would, using more resources for prediction and processing. As a result, time can feel like it “slows down” in situations of high stress.

These examples show the complexity and nuance in how time is perceived. Dopamine and norepinephrine are the primary modulators, but every neurotransmitter plays a role in how we perceive reality and therefore how we perceive time. Calculating how species perceive reality and time is incredibly complex and nonlinear. It is not just a matter of basic sensory functions but a huge neurological network with high variance governing how stimulus is perceived and integrated, how brain energy is allocated, how neurotransmitters are mapped, and the nuance with which time perception is impacted.

\section{Formal mathematical framework}

This section proposes a conceptual model of biological time perception in standard mathematical notation. The framework functions in two layers: 1) an instantaneous layer that captures how much conscious content fills a given moment and 2) a retroactive layer that captures the felt speeding-up or slowing-down of time by comparing the present moment against a trailing memory baseline.

\medskip

The table below lists every symbol used throughout:

\begin{center}
	\renewcommand{\arraystretch}{1.3}
	\begin{tabular}{@{}lll@{}}
		\toprule
		Symbol & Name & Range \\
		\midrule
		$t$ & continuous time & $t \ge 0$ \\
		$C$ & perceptual ceiling (species constant) & $C > 0$ \\
		$x_e(t)$ & external sensory input & $[0,\,C]$ \\
		$x_s(t)$ & somatic (body) input & $[0,\,C]$ \\
		$x_m(t)$ & mental/thought input & $[0,\,C]$ \\
		$\varphi_i(t)$ & familiarity filter for channel $i$ & $\ge 0$ \\
		$E_i(t)$ & entropy (disorder) of the stimulus in channel $i$ & $\ge 0$ \\
		$\kappa$ & baseline clock-monitoring intensity & $> 0$ \\
		$K(t)$ & clock-checking load & $[0,\,C]$ \\
		$D(t)$ & subjective temporal density & $[0,\,C]$ \\
		$W$ & memory window scale & $> 0$ \\
		$\Psi(t)$ & trailing memory contrast index & $(-C,\,C)$ \\
		\bottomrule
	\end{tabular}
\end{center}

The subscript $i \in \{e,\,s,\,m\}$ labels the three input channels: external sensory, somatic, and mental.

\subsection{Layer 1 -- Instantaneous temporal density}

At every instant $t$, the brain takes in three parallel data streams:
\[
\mathbf{x}(t) \;=\; \bigl(\, x_e(t),\;\; x_s(t),\;\; x_m(t) \,\bigr).
\]
Each $x_i(t)$ sits between $0$ and $C$ and measures the raw information load coming in through that channel.

Two neural processes act as expanders or contractors shaping how that raw data load is processed. They work in opposite directions:

Familiarity ($\varphi_i$) compresses. The more familiar a brain is with a stimulus, the less processing it needs. A familiar commute barely registers; a brand-new city street demands full attention.

Entropy ($E_i$) amplifies. Entropy here means disorder in the environment \citep{Hechler2023energy}. When a scene is orderly and predictable ($E_i \approx 0$), the brain handles it easily. But when disorder spikes (e.g., a car crash with fire, smoke, and flying debris), the brain ramps up predictive processing to keep the organism safe. The effective load goes up \citep{Tang2025less}.

The effective load from channel $i$ after both filters is:
\begin{equation}
	L_i(t) \;=\; \frac{x_i(t)\;\bigl(1 + E_i(t)\bigr)}{1 + \varphi_i(t)}\,.
\end{equation}

Familiarity sits in the denominator (compresses), entropy sits in the numerator (amplifies). When both are zero -- a novel but orderly stimulus -- the full raw load passes through unchanged: $L_i = x_i$. In a familiar, calm setting ($\varphi_i$ high, $E_i$ low), the denominator dominates and the brain barely registers it. In a chaotic, unfamiliar crisis ($\varphi_i$ low, $E_i$ high), the load can exceed the raw input because the brain is working overtime to predict the unfolding disorder.

The total task load across all three channels is simply their sum:
\begin{equation}
	\Lambda(t) \;=\; L_e(t) + L_s(t) + L_m(t).
\end{equation}

Clock checking is an attentional process that competes with those three channels for space under the perceptual ceiling. The brain cannot use familiarity to compress its own internal timing pulses, so the clock-checking load $K$ depends only on how much ceiling capacity the task load leaves unused \citep{zakay1996attention}:
\begin{equation}
	K(t) \;=\; \kappa\;\left[\,1 - \frac{\Lambda(t)}{C}\,\right]^{+}
\end{equation}

where $[\,\cdot\,]^{+} = \max(\,\cdot\,,\;0\,)$. If the task load is dense and engaging ($\Lambda \to C$), there is no room left for clock monitoring and the gate slams shut ($K \to 0$). If the task load is low and understimulating ($\Lambda \to 0$), the gate swings wide open ($K \to \kappa$) and the organism floods its remaining capacity with clock awareness.

Putting it together, the instantaneous subjective temporal density is the total volume of conscious content at that moment, capped at the species ceiling:
\begin{equation}
	D(t) \;=\; \min\!\bigl(\,\Lambda(t) + K(t),\;\; C\,\bigr).
\end{equation}

A normalised version on a 0-to-1 scale is $\;d(t) = D(t)/C$.

\subsection{Layer 2 -- Retroactive mechanism (the temporal window)}

Once the present-moment density $D(t)$ is derived from the core equation, the brain evaluates whether time is speeding up or slowing down by comparing that density against a trailing average over a memory window of depth $W$:
\begin{equation}
	\bar{D}_W(t) \;=\; \frac{1}{W}\int_{t-W}^{\,t} D(\tau)\;\mathrm{d}\tau\,.
\end{equation}

This is the memory baseline -- the average density the brain has experienced over the past $W$ time units. The parameter $W$ (the memory window scale) is the independent variable that sets how far back the comparison reaches.

The trailing memory contrast index is then the signed difference between the current density and that baseline:
\begin{equation}
	\Psi(t,\,W) \;=\; D(t) \;-\; \bar{D}_W(t).
\end{equation}

\begin{center}
	\renewcommand{\arraystretch}{1.3}
	\begin{tabular}{@{}cl@{}}
		\toprule
		Sign of $\Psi$ & What it feels like \\
		\midrule
		$\Psi > 0$ & Time feels like it is slowing down (richer moment than baseline) \\
		$\Psi < 0$ & Time feels like it is speeding up (sparser moment than baseline) \\
		$\Psi \approx 0$ & No noticeable change in the rate of time \\
		\bottomrule
	\end{tabular}
\end{center}

\medskip

To illustrate, consider the dark-shed thought experiment from \S~6.1. Inside the shed, external input drops to near zero and familiarity builds up, so $\Lambda$ stays low and clock checking ($K$) is high. The density $D$ settles to a high baseline steady state driven by primarily internal data amplified by clock awareness. If someone steps outside after just 10 seconds, the memory window $W$ is small and the baseline $\bar{D}_W$ has barely shifted, so $\Psi$ registers only a mild jolt. But if someone sits in that shed for a full year, the baseline gets increased by twelve months of dense, internally-driven data. The moment they step outside, $\Psi$ drops sharply -- cognitive shock. And normalisation is slow, because the running average needs enough outdoor data to catch up.

Long-term sensory deprivation triggers homeostatic plasticity, restructuring the brain to hyper-process internal somatic signals and generate multi-sensory hallucinations. This occurs as the damaged hippocampus and altered dopaminergic pathways degrade the brain's ability to filter background noise, causing the brain to manufacture its own high-intensity internal data stream, which increases overall net stimulus \citep{marazziti2026sensory}.

\subsection{Cross-species extension}

The perceptual ceiling $C$ is a species-level constant. For typical mammals it can be roughly approximated through the visual system:
\begin{equation}
	C \;\approx\; \alpha_1 \cdot \mathrm{CFF}
	\;+\; \alpha_2 \cdot \mathrm{FOV}
	\;+\; \alpha_3 \cdot R
	\;+\; \sum_{j}\,\omega_j\,S_j\,,
\end{equation}

where CFF is the temporal resolution (like FPS), FOV is the binocular field of view, $R$ is the ocular spatial resolution (acuity), $S_j$ is the capacity of any additional sensory modality $j$ (echolocation in bats, electroreception in sharks, etc.), $\omega_j$ is a species-specific weight reflecting how heavily the organism relies on that modality, and $\alpha_1, \alpha_2, \alpha_3$ are scaling constants to be fitted from empirical data. Different species allocate bandwidth differently across their senses, so the weights matter.

One important prerequisite: the entire framework assumes the organism has a clear boundary between self and stimulus. Without that distinction, the comparison between present and past cannot happen, and there is no individual to perceive time (e.g., worms). Formally, $\Psi$ is undefined for such organisms.

\subsection{Complete equation set}

Given: species ceiling $C$, baseline clock intensity $\kappa$, memory window $W$, input streams $x_e, x_s, x_m$ with their familiarity factors $\varphi_i$ and entropy levels $E_i$.

\begin{align}
	L_i(t) &= \frac{x_i(t)\,(1+E_i(t))}{1+\varphi_i(t)},
	\quad i\in\{e,s,m\}
	\tag{1}\\[6pt]
	\Lambda(t) &= L_e(t)+L_s(t)+L_m(t)
	\tag{2}\\[6pt]
	K(t) &= \kappa\left[1-\frac{\Lambda(t)}{C}\right]^{+}
	\tag{3}\\[6pt]
	D(t) &= \min\!\bigl(\Lambda(t)+K(t),\; C\bigr)
	\tag{4}\\[6pt]
	\bar{D}_W(t) &= \frac{1}{W}\int_{t-W}^{t}D(\tau)\,\mathrm{d}\tau
	\tag{5}\\[6pt]
	\Psi(t,W) &= D(t) - \bar{D}_W(t)
	\tag{6}
\end{align}

\medskip
$\Psi > 0$: time feels slower. \quad
$\Psi < 0$: time feels faster. \quad
$\Psi \approx 0$: no perceived change.

\subsection{Analysis}

Entropy ($E_i$) in this model refers to physical/environmental disorder, not the information-theoretic quantity. A car crash, an explosion, a sudden storm -- these are high-entropy events where the environment is rapidly expanding towards disorder. The brain must track many more unpredictable elements (debris, fire, smoke, shifting trajectories), which forces it to process more data to predict what comes next. That is why entropy appears as a multiplier in the numerator: more disorder means more work for the brain.

This framework incorporates five primary inputs for the perceptual ceiling: somatic, external sensory, mental thought, predictive processing, and subcortical. The equations above model the first three as explicit channels. Predictive processing is captured through the interplay of the two filters -- familiarity handles learned priors, entropy handles the degree of environmental disorder the brain must predict through. Subcortical contributions (autonomic arousal, brainstem alertness, etc.) can be added as an extra term in $\Lambda(t)$ once empirical data become available.

All variables are written as continuous functions of time. For computational simulation, the integral in equation (5) can be replaced by a discrete sum over time steps $\Delta t$, and $D(t)$ can be computed at each step from equations (1) through (4).

Finally, while the conceptual framework of the preceding equation is viable, the specific notation requires further peer review and refinement.

\section{Thought experiment}

The following cases illustrate the mechanisms described in \S~4. They are not independent proofs.

The world line refers to the frozen coordinates of a life from beginning to end within the static spacetime block. As an example, consider a world line that runs for exactly 100 years. At the one-quarter mark, the person is 25 years old; at halfway, he is 50; at three-quarters, he is 75; and so on. In this thought experiment, there is a hypothetical moving present, a play head that moves along the worldine representing ``now.''

If we were to hypothetically cut a world line in half at 50 years, what exactly does that mean? We can think of it like cutting a video. Both parts still exist but instead of returning to the very beginning, the hypothetical playhead returns to wherever it was cut and repeats from there. So, would he still experience the illusion of a full continuous lifetime? Yes, because nothing within the actual world line was changed or removed. The conscious being within that world line still believes he's lived a completely continuous life as he remembers the other half of his world line and is convinced of his own continuity.

Now, what if we cut it again at 25 years? The same principle applies. The subject then remembers the first 25 years. The same is true if we cut at 50, 75, and so on. Since they are being divided at the same interval, they each remain at the same point along their world line.

However, while we are dividing evenly in this experiment for simplicity, the scale of the division is arbitrary. It doesn't need to be even. The same principles apply, regardless of the interval size.

If we continue cutting the world line, eventually we reach the millisecond scale where perception is divided into micro units. This is where frames come in. For exampe, a few thousand milliseconds might equate to 5 perceptual frames. Frames are the integrated chunks of experience our brain creates to stitch sensory input into a coherent experience of conscious snapshots.

As in the previous example, because frames only contain memory data from prior frames, the frame where the observer is located in their micro worldline is completely convinced of its own continuity.

If we continue to narrow down and isolate these units, there are no longer any hypothetical ``off'' frames. The playhead no longer moves as there is no longer a unique ``now'' and each frame is now ``on.'' Each static frame within the worldline is now actively conscious. These are the mechanics of concsciousness within the block. 

If you want to claim it's somehow different than prior examples, you introduce a ghost mechanism by implying that our perception of time operates outside the very constraint it must operate within, which is perception. You are claiming that this unit is not sufficient to give rise to a coherent experience, disregarding the field of neuroscience, and denying experience itself.

Frames don't sit there and wait around until the other frames are not active to finally become conscious. A frame within the block is forever a conscious snapshot completely convinced of its own continuity regardless of its location on the world line. All it knows is what came before it. As you read this, you believe yesterday was yesterday and tomorrow is about to happen because within the block, your neurological state is frozen. It makes absolutely zero difference, whether as in the prior example, four frames were inactive or if every frame is active. There is not some magical force outside these frames responsible for the perception of time passing, no magical fairy outside the spacetime block that drifts along the world line permitting perception one frame at a time.

As another example, consider a subject observing a table. Is it possible to know which side of the table is the beginning and which is the end? No. Like the world line within the block, there is no inherent visual directionality. The subject might think, “I know life began at birth and ended at death,” but that's due to prior knowledge about human beings. An entity who doesn't understand how humans mature would see it exactly how a subject observes a blank table—just an orderless, directionless structure.

We can attempt to measure this object from the outside, but that provides no inherent causal structure. We are adding external mechanisms which have no impact on its inherent mechanics. This is the approach taken by models like C theory. They attempt to provide a causal structure by drawing on it. This is what I refer to as a ghost mechanism.

Similarly, using entropy to explain temporal asymmetry is a fallacy and is unnecessary.

This hypothetical table, for example, has inherent causal structure, so it is built into the table itself rather than trying to force a causal structure from the outside. So, A remembers B, B remembers C, C remembers D, and so on.

\section{Limits and implications}

This paper is not an attempt to prove that time is not real, or that physics is incorrect about the structure of spacetime. Instead, the paper has demonstrated that the felt experience of time passing does not require motion of any sort. The continuous, smooth experience of time can instead be generated entirely by the manner in which our brain integrates stimulus, the mechanics of the block and the static brain contents within it, and the fixed, one-way dependency between these perceptual frames.

The biological and neurological data used to test that mechanism and explain how humans and other species experience time is drawn from existing data. The more speculative claims regarding dopamine and norepinephrine signaling are this paper’s interpretive extension of existing data and the neurological complexity governing our perception of time.

\section{Conclusion}

If the universe is a static block and all moments exist simultaneously, how do we experience the smooth passage of time? Rather than introducing external variables, the answer lies within neurology itself. The mechanics of consciousness—and how our brain integrates data—provide a simple and inherent explanation for the fluid experience of time and why we can only remember the past. Nothing needs to sweep across the block or be imported from the outside.

\bibliographystyle{apalike}
\bibliography{references}

\end{document}
